% !TITLE 泊秦淮
% !AUTHOR 杜牧
% !DYNASTY 唐
% !GENRE 七言绝句
% !ORDER 56

\begin{poem}
烟笼寒水月笼沙\textsuperscript{1}，夜泊秦淮\textsuperscript{2}近酒家。
商女\textsuperscript{3}不知亡国恨，隔江犹唱后庭花\textsuperscript{4}。
\end{poem}

\begin{annotations}
\notetext{1}{烟笼寒水月笼沙：烟霭笼罩着寒水，月光笼罩着沙滩。两个“笼”字叠用，为全诗罩上一层朦胧的色调。}
\notetext{2}{秦淮：秦淮河，流经南京，六朝以来最繁华的风月之地，名门士族宴饮游乐之处。}
\notetext{3}{商女：以歌唱为业的女子。歌女不懂亡国之痛——但真正“不知亡国恨”的，是那些在酒家中寻欢作乐的达官贵人。}
\notetext{4}{后庭花：即《玉树后庭花》，南朝陈后主所作的艳曲，被认为是亡国之音。隔江还能听到歌女在唱这支曲子——晚唐的达官仍在重蹈六朝的覆辙。}
\end{annotations}

\begin{appreciation}
这是晚唐最著名的讽刺诗之一，二十八个字，把一个王朝的病写完了。杜牧生活的时代，大唐已经走到最后几十年，他自己仕途也不顺，在幕府间辗转。

烟笼寒水月笼沙——两个笼字：烟罩着水，月罩着沙。不是清朗的夜色，是迷蒙的、什么都看不真切的夜色——像这个王朝的局势，也像酒家灯火里那些人的醉眼。

夜泊秦淮近酒家——秦淮河，流过金陵，六朝以来最繁华的风月之地。泊，停船。近，只是靠近，没有进去。杜牧是隔着一层听的。

商女不知亡国恨，隔江犹唱后庭花——《玉树后庭花》，南朝陈后主作的艳曲，他就在金陵，唱着这支曲子唱到亡国。两百多年过去，同一片水面上，歌女又唱起来了。犹字最重：还在唱。历史不是重演，历史根本就没走。

这里有个转折要读懂。歌女唱什么，由听众定——真正不知亡国恨的，不是唱歌的人，是那些点歌的人：酒家里的达官贵人，醉生梦死的座上宾。杜牧不骂他们，只淡淡写一句隔江犹唱，比骂一万句都重。

还记得《生年不满百》吗？两汉的人说昼短苦夜长，何不秉烛游——乱世里人生苦短，及时行乐，是无奈的聪明。可同一个秉烛游，搬到晚唐的秦淮河上，就成了亡国前的麻醉。同样是灯红酒绿，时代替它分出了高下。这本书后面，王安石站在同一座城，在《桂枝香·金陵怀古》里写下六朝旧事随流水——六朝的旧事随水流走了，歌却没有停。

你将来去南京，晚上路过秦淮河，如果听见有人唱歌，可以想想这首诗。那条河，六朝唱到唐，唐唱到今天。
\end{appreciation}
